\documentclass[uplatex,dvipdfmx]{jsarticle}
\usepackage{amsmath,amssymb,amsthm}
\usepackage{thmtools,thm-restate}
\usepackage{enumitem}
\usepackage{seqsplit}
\usepackage[dvipdfmx,unicode]{hyperref}
\hypersetup{
  pdftitle={SOS Factorization and Markov--Lukacs Certificates for Real Symmetric Univariate Matrix Polynomials},
  pdfauthor={Mocho Go}
}

\newcommand{\Poly}{\mathbb{R}[X]}
\newcommand{\Frac}{\mathbb{R}(X)}
\newcommand{\RationalHahn}{\operatorname{RHahn}}
\newcommand{\Hahn}{\mathbb{R}((t^{\mathbb{Q}}))}
\newcommand{\lc}{\operatorname{lc}}
\newcommand{\ord}{\operatorname{ord}}
\newcommand{\Mat}{\operatorname{Mat}}
\newcommand{\psd}{\succeq 0}
\newcommand{\NNwd}{\operatorname{NNwd}}
\newcommand{\lean}[1]{\texttt{\expandafter\seqsplit\expandafter{\detokenize{#1}}}}

\newcommand{\restateprefix}{}
\newtheoremstyle{matrixsosdefinition}
  {3pt}{3pt}{\normalfont}{}{\bfseries}{.}{.5em}
  {\thmname{#1}\thmnumber{ #2}\thmnote{ (\restateprefix#3)}}

\theoremstyle{matrixsosdefinition}
\newtheorem{theorem}{Theorem}
\newtheorem{lemma}[theorem]{Lemma}
\newtheorem{definition}[theorem]{Definition}
\newtheorem{corollary}[theorem]{Corollary}

\title{SOS Factorization and Markov--Lukacs Certificates for Real Symmetric Univariate Matrix Polynomials}
\author{Mocho Go}
\date{}

\begin{document}
\maketitle

\section{Overview and Main Theorems}

The central result formalized here is a bounded full-line sums-of-squares
factorization theorem for real symmetric univariate matrix polynomials.  As
corollaries, the formalization also gives half-line Markov--Lukacs-type SOS certificates,
finite-interval Markov--Lukacs-type SOS certificates, and scalar and constant
matrix special cases.  The proof uses a Hahn series field with rational
exponents as an auxiliary ordered field; the final statements are statements
about real polynomial matrices.  The full-line theorem follows the mathematical
theme of Hanselka--Sinn~\cite{HanselkaSinn}.

Here SOS abbreviates ``sum of squares''.  In the scalar case this means a finite
sum of polynomial squares; in the matrix case it means the corresponding
rectangular factorization
\[
  S=R^{\mathsf T}R
\]
by a polynomial matrix \(R\).

The Markov--Lukacs-type statements are the half-line and interval forms that
express nonnegativity using SOS terms weighted by polynomials nonnegative on the
domain.  In this project the weights include \(X\), \(X-a\), \(b-X\), and
\((X-a)(b-X)\).  For the classical scalar background, see
Powers--Reznick~\cite{PowersReznick}.  Related work on matrix-polynomial
positivity and SOS representations includes
Blekherman--Plaumann--Sinn--Vinzant~\cite{BPSV} and
Klep--Schweighofer~\cite{KlepSchweighofer}.

For a real symmetric matrix \(A\), \(A\psd\) means
\(v^{\mathsf T}Av\ge0\) for every column vector \(v\).  Thus \(P(x)\psd\) for a
polynomial matrix \(P\) means that the real matrix obtained by evaluating each
entry at \(x\) is positive semidefinite.

\begin{definition}[SOS matrix polynomials and Gram representations]\label{def:sos-gram-en}
  Let \(m\in\mathbb{N}\).
  A matrix polynomial \(S\in\Mat_m(\Poly)\) is SOS if there is a polynomial
  matrix \(R\in\Mat_{\ell\times m}(\Poly)\), for some \(\ell\in\mathbb{N}\),
  such that
	  \[
	    S=R^{\mathsf T}R.
	  \]
	  We call such an \(R\) an SOS factor of \(S\).

	  For \(r,m\in\mathbb{N}\), define
	  \(Q_{r,m}\in\Mat_{(r+1)m\times m}(\Poly)\) as follows.  Its rows are
	  indexed by pairs \((k,a)\), \(0\le k\le r,\ 1\le a\le m\), its columns are
	  indexed by \(1\le b\le m\), and its entries are
  \[
    (Q_{r,m})_{(k,a),b}=
    \begin{cases}
      X^k, & a=b,\\
	      0, & a\ne b.
	    \end{cases}
	  \]
	  For any real matrix \(Y\in\Mat_{(r+1)m}(\mathbb{R})\), the product
	  \(Q_{r,m}^{\mathsf T}YQ_{r,m}\) is an \(m\times m\) polynomial matrix; if
	  \(Y\) is symmetric, then this polynomial matrix is symmetric.
	  A bounded SOS factorization of \(S\) with factor degree at most \(e\) is
	  equivalent to a Gram representation
  \[
    S=Q_{e,m}^{\mathsf T}YQ_{e,m}
  \]
	  with a real positive semidefinite matrix \(Y\).  This is obtained by expanding
  the entries of the SOS factor in the monomial basis \(1,X,\ldots,X^e\) and
  collecting the coefficients into a real matrix.  Conversely, the spectral
  theorem gives \(Y=C^{\mathsf T}C\) for a real positive semidefinite \(Y\);
  then \(R=CQ_{e,m}\) is an SOS factor of degree at most \(e\).
\end{definition}

\begin{restatable}[Full-line matrix SOS]{theorem}{fullLineMatrixSOS}
  Let \(m,d\in\mathbb{N}\). Let \(P\in\Mat_m(\Poly)\) be symmetric, and assume
  every entry of \(P\) has degree at most \(2d\). Then the following are
  equivalent:
  \begin{enumerate}[label=(\arabic*)]
    \item \(P(x)\psd\) for every \(x\in\mathbb{R}\).
    \item There is \(R\in\Mat_{(m+1)\times m}(\Poly)\), whose entries all have
      degree at most \(d\), such that
      \[
        P=R^{\mathsf T}R.
      \]
    \item There is a real positive semidefinite matrix
      \(Y\in\Mat_{(d+1)m}(\mathbb{R})\) such that
      \[
        P=Q_{d,m}^{\mathsf T}YQ_{d,m}.
      \]
  \end{enumerate}
\end{restatable}

\begin{restatable}[Half-line Markov--Lukacs certificate]{theorem}{halfLineMarkovLukacs}
  Let \(m,d\in\mathbb{N}\), let \(M\in\Mat_m(\Poly)\) be symmetric, and let
  \(a\in\mathbb{R}\).
  Assume every entry of \(M\) has degree at most \(d\).  Put
  \[
    r(d)=\left\lfloor\frac{d}{2}\right\rfloor,\qquad
    s(d)=\left\lfloor\frac{\max(d-1,0)}{2}\right\rfloor.
  \]
  Then the following are equivalent:
  \begin{enumerate}[label=(\arabic*)]
    \item \(M(x)\psd\) for every \(x\ge a\).
    \item There are \(A,B\in\Mat_{(m+1)\times m}(\Poly)\) such that
      \[
        M=A^{\mathsf T}A+(X-a)B^{\mathsf T}B,
      \]
      with entries of \(A\) of degree at most \(r(d)\), and entries of \(B\) of
      degree at most \(s(d)\).
    \item There are real positive semidefinite matrices
      \[
        Y_1\in\Mat_{(r(d)+1)m}(\mathbb{R}),\qquad
        Y_2\in\Mat_{(s(d)+1)m}(\mathbb{R})
      \]
      such that
      \[
        M =
        Q_{r(d),m}^{\mathsf T}Y_1Q_{r(d),m}
        +
        (X-a)\,Q_{s(d),m}^{\mathsf T}Y_2Q_{s(d),m}.
      \]
  \end{enumerate}
\end{restatable}

\begin{restatable}[Finite-interval Markov--Lukacs certificate]{theorem}{finiteIntervalMarkovLukacs}
  Let \(m,d\in\mathbb{N}\), let \(M\in\Mat_m(\Poly)\) be symmetric, and let
  \(a,b\in\mathbb{R}\) with \(a<b\).
  Put \(d_-=\max(d-1,0)\).
  If the entries of \(M\) have degree at most \(2d\), then the following are
  equivalent:
  \begin{enumerate}[label=(E\arabic*)]
    \item \(M(x)\psd\) for every \(x\in[a,b]\).
    \item There are \(A,B\in\Mat_{(m+1)\times m}(\Poly)\) such that
      \[
        M=A^{\mathsf T}A+(X-a)(b-X)B^{\mathsf T}B,
      \]
      with entries of \(A\) of degree at most \(d\), and entries of \(B\) of
      degree at most \(d_-\).
    \item There are real positive semidefinite matrices
      \[
        Y_0\in\Mat_{(d+1)m}(\mathbb{R}),\qquad
        Y_1\in\Mat_{(d_-+1)m}(\mathbb{R})
      \]
      such that
      \[
        M =
        Q_{d,m}^{\mathsf T}Y_0Q_{d,m}
        +(X-a)(b-X)\,Q_{d_-,m}^{\mathsf T}Y_1Q_{d_-,m}.
      \]
  \end{enumerate}
  If the entries of \(M\) have degree at most \(2d+1\), then the following are
  equivalent:
  \begin{enumerate}[label=(O\arabic*)]
    \item \(M(x)\psd\) for every \(x\in[a,b]\).
    \item There are \(A,B\in\Mat_{(m+1)\times m}(\Poly)\) such that
      \[
        M=(X-a)A^{\mathsf T}A+(b-X)B^{\mathsf T}B,
      \]
      with entries of \(A\) and \(B\) of degree at most \(d\).
    \item There are real positive semidefinite matrices
      \(Y_0,Y_1\in\Mat_{(d+1)m}(\mathbb{R})\) such that
      \[
        M =
        (X-a)\,Q_{d,m}^{\mathsf T}Y_0Q_{d,m}
        +(b-X)\,Q_{d,m}^{\mathsf T}Y_1Q_{d,m}.
      \]
  \end{enumerate}
\end{restatable}

\begin{corollary}[Scalar and constant cases]
  Let \(d\in\mathbb{N}\), and let \(p\in\mathbb{R}[X]\) have degree at most
  \(2d\).  Then \(p\) is nonnegative on the real line if and only if
  \(p=q_1^2+q_2^2\) for some real polynomials \(q_1,q_2\) of degree at most
  \(d\).
  Let \(m\in\mathbb{N}\), and let \(A\in\Mat_m(\mathbb{R})\).  Then \(A\) is
  positive semidefinite if and only if \(A=R^TR\) for some
  \(R\in\Mat_{(m+1)\times m}(\mathbb{R})\).
\end{corollary}

The proof below uses the following rational-function lemma.  Its proof is the
part that uses the Hahn field, Tsen's theorem, and Pfister forms, and is given
in Appendices~\ref{sec:hahn-start-en}--\ref{sec:ternary-local-global-en}.

First define the notation used in the statement.  For \(r\in\Frac\), write
\(\NNwd(r)\) if \(r\) is nonnegative at every real point where it is defined:
for a representation \(r=f/g\) with \(g\ne0\),
\[
  g(x)\ne0 \quad\Longrightarrow\quad \frac{f(x)}{g(x)}\ge0
  \qquad (x\in\mathbb{R}).
\]
This condition is independent of the chosen representation.

\begin{restatable}[Binary sum-of-squares representation]{lemma}{binarySumTwoSquares}
\label{lem:binary-sum-two-squares-en}
  If \(a,b\in\Frac^\times\) satisfy \(\NNwd(a)\) and \(\NNwd(b)\), then
  there exist \(u,v\in\Frac\) such that
  \[
    1=au^2+bv^2.
  \]
\end{restatable}

We first prove the matrix part assuming
Lemma~\ref{lem:binary-sum-two-squares-en}.  The proof of that lemma is given in
the appendix.

\section{Full-Line Matrix Factorization}\label{sec:matrix-factorization-en}

\begin{lemma}\label{lem:rational-square-factorization-en}
  Let \(n\in\mathbb{N}\). Let \(A\in\Mat_n(\Poly)\) be positive semidefinite at
  every real point, and suppose \(\det A\ne0\) and \(\det A\) is a square in
  \(\Frac\). Then
  \[
    A=S^{\mathsf T}S
  \]
  for some \(S\in\Mat_n(\Frac)\).
\end{lemma}

\begin{proof}
  If \(n=0\), the unique empty matrix is the required factor.  Assume \(n>0\).
  Put \(F=\Frac\). Diagonalize the symmetric matrix \(A\) over \(F\). Thus
  there are \(C\in\operatorname{GL}_n(F)\) and
  \(d_1,\ldots,d_n\in F^\times\) such that
  \[
    C^{\mathsf T}AC=\operatorname{diag}(d_1,\ldots,d_n).
  \]
  Positivity at real points shows that each \(d_i\) is nonnegative where
  defined: outside the finite set of poles of \(C\), evaluate the congruence;
  at any remaining point where \(d_i\) is defined, use continuity of \(d_i\)
  and density of the complement of that finite set. Since \(\det A\) is a square, write \(\det A=r^2\) with \(r\in F\).
  Then
  \[
    d_1\cdots d_n
    =\det(C^{\mathsf T}AC)
    =\det(C)^2\det(A)
    =(\det(C)r)^2,
  \]
  so the product of the diagonal coefficients is also a square in \(F\).

  The key two-dimensional move is the following. If \(a,b\in\Frac^\times\) are
  nonnegative where defined, Lemma~\ref{lem:binary-sum-two-squares-en} gives
  \(1=au^2+bv^2\). With
  \[
    G=
    \begin{pmatrix}u&-bv\\ v&au\end{pmatrix}\in\operatorname{GL}_2(F),
  \]
  a direct calculation gives
  \[
    G^{\mathsf T}
    \begin{pmatrix}a&0\\0&b\end{pmatrix}
    G
    =
    \begin{pmatrix}1&0\\0&ab\end{pmatrix}.
  \]
  Here \(\det G=a u^2+b v^2=1\).  Thus this congruence transformation
  extracts a square entry \(1\) from the two diagonal entries \(a,b\) and
  replaces the remaining information by their product \(ab\). Applying this
  operation successively gives a matrix \(E\in\operatorname{GL}_n(F)\) such that
  \[
    E^{\mathsf T}\operatorname{diag}(d_1,\ldots,d_n)E
    =
    \operatorname{diag}(1,\ldots,1,d_1\cdots d_n).
  \]
  Since the last entry is a square, say \(d_1\cdots d_n=(r')^2\) with
  \(r'\in F\), the final diagonal form is \(D^{\mathsf T}D\), where
  \[
    D=\operatorname{diag}(1,\ldots,1,r')\in\Mat_n(F).
  \]
  If \(C^{\mathsf T}AC=\operatorname{diag}(d_1,\ldots,d_n)\) is the initial
  diagonalization and \(E\) is the product of the two-dimensional changes, then
  \(E^{\mathsf T}C^{\mathsf T}ACE=D^{\mathsf T}D\). Hence
  \(A=(D(CE)^{-1})^{\mathsf T}(D(CE)^{-1})\).
\end{proof}

\begin{lemma}\label{lem:clear-common-denominator-en}
  Suppose \(A\in\Mat_n(\Poly)\) and \(S\in\Mat_n(\Frac)\) satisfy
  \(A=S^{\mathsf T}S\). Then there are a nonzero polynomial \(q\in\Poly\) and
  \(T\in\Mat_n(\Poly)\) such that
  \[
    q^2A=T^{\mathsf T}T.
  \]
\end{lemma}

\begin{proof}
  Write \(A=S^{\mathsf T}S\) over \(\Frac\). Let \(q\) be a common denominator
  of all entries of \(S\), and put \(T=qS\). Then
  \(T^{\mathsf T}T=q^2A\).
\end{proof}

\begin{lemma}[Real-root factor descent]\label{lem:real-root-factor-descent-en}
  Let \(p\in\Poly\) be irreducible, and suppose \(p(\alpha)=0\) for some
  \(\alpha\in\mathbb{R}\).  If \(A\in\Mat_n(\Poly)\) and
  \(T\in\Mat_n(\Poly)\) satisfy
  \[
    p^2A=T^{\mathsf T}T,
  \]
  then there is \(U\in\Mat_n(\Poly)\) such that
  \[
    A=U^{\mathsf T}U.
  \]
\end{lemma}

\begin{proof}
  Comparing diagonal entries gives
  \[
    p^2A_{ii}=\sum_k T_{ki}^2.
  \]
  Evaluating at \(X=\alpha\), the left hand side is zero, so
  \[
    \sum_k T_{ki}(\alpha)^2=0.
  \]
  This is a sum of squares of real numbers, hence \(T_{ki}(\alpha)=0\) for all
  \(k,i\).  Since an irreducible real polynomial with root \(\alpha\) is a
  nonzero scalar multiple of \(X-\alpha\), every \(T_{ki}\) is divisible by
  \(p\).  Write \(T=pU\), with
  \(U\in\Mat_n(\Poly)\).  Substitution gives
  \(p^2A=p^2U^{\mathsf T}U\), and cancellation of \(p^2\) gives
  \(A=U^{\mathsf T}U\).
\end{proof}

\begin{lemma}[Two-row block construction]\label{lem:two-row-block-construction-en}
  Let \(q,n\in\mathbb{N}\), and let \(N=2q\) or \(N=2q+1\).
  Let \(a,b\in\mathbb{R}\) with \(b\ne0\), and put
  \[
    p=(X-a)^2+b^2\in\Poly.
  \]
  Let \(T\in\Mat_{N\times n}(\Poly)\), and suppose that \(p\) divides every
  entry of \(T^{\mathsf T}T\).
  Then there are a real orthogonal matrix \(O\in\Mat_N(\mathbb{R})\) and
  matrices \(H\in\Mat_N(\Poly)\) and \(M\in\Mat_{N\times n}(\Poly)\) with the
  following properties.  If \(N=2q\), then
  \[
    OT=HM,\qquad H^{\mathsf T}H=pI_N.
  \]
  If \(N=2q+1\), then
  \[
    OT=HM,\qquad
    H^{\mathsf T}H=\operatorname{diag}(pI_{2q},p^2).
  \]
\end{lemma}

\begin{proof}
  Put \(z=a-ib\in\mathbb{C}\).  For a polynomial \(f\in\Poly\), divisibility by \(p\) is
  equivalent to \(f(z)=0\).

  From the divisibility assumption, \(T(z)^{\mathsf T}T(z)=0\).  Thus the
  subspace \(L\subset\mathbb{C}^{N}\) generated by the columns of \(T(z)\) is
  totally isotropic for the bilinear form \(x^{\mathsf T}y\).  Choose a
  Hermitian orthonormal basis \(w_1,\ldots,w_r\) of \(L\).  Write
  \(w_j=a_j+ib_j\), with \(a_j,b_j\in\mathbb{R}^{N}\). Then
  \[
    \sqrt2\,\Re w_1,\ \sqrt2\,\Im w_1,\ldots,
    \sqrt2\,\Re w_r,\ \sqrt2\,\Im w_r
  \]
  is an orthonormal system in \(\mathbb{R}^{N}\).  Hence \(2r\le N\), and
  \(r\le\lfloor N/2\rfloor=q\).
  Extend these vectors to an orthonormal basis \(e_1,\ldots,e_N\).  Let
  \(O\in\Mat_N(\mathbb{R})\) be the real orthogonal matrix whose
  \(a\)-th row is \(e_a^{\mathsf T}\).

  If \(y\) is any column of \(T(z)\), then \(y\in L\), so
  \(y=\sum_j\mu_jw_j\).  For \(1\le j\le r\), the two coordinates of \(y\)
  measured by \(e_{2j-1},e_{2j}\) are
  \[
    \left(\sqrt2\,a_j^{\mathsf T}y,\sqrt2\,b_j^{\mathsf T}y\right)
    =\frac{\mu_j}{\sqrt2}(1,i).
  \]
  For the remaining basis vectors, the corresponding rows of \(OT(z)\) are
  zero.

  Put \(B=\begin{pmatrix}X-a&-b\\ b&X-a\end{pmatrix}\).  Then
  \(B^{\mathsf T}B=BB^{\mathsf T}=pI_2\).  Also
  \(B(z)^{\mathsf T}(1,i)^{\mathsf T}=0\).  For each \(1\le j\le q\), let
  \((r_1,r_2)\) be the two-row pair of \(OT\) in rows \(2j-1,2j\).  The matrix
  \[
    B^{\mathsf T}\begin{pmatrix}r_1\\r_2\end{pmatrix}
  \]
  vanishes at \(z\).  Since \(p\mid f\Longleftrightarrow f(z)=0\), all entries
  are divisible by \(p\).  Define row vectors
  \(m_1,m_2\in\Mat_{1\times n}(\Poly)\) by
  \[
    B^{\mathsf T}\begin{pmatrix}r_1\\r_2\end{pmatrix}
      =
    p\begin{pmatrix}m_1\\m_2\end{pmatrix}.
  \]
  Multiplying by \(B\) and using \(BB^{\mathsf T}=pI_2\), then cancelling \(p\),
  gives
  \[
    \begin{pmatrix}r_1\\r_2\end{pmatrix}
      =
    B\begin{pmatrix}m_1\\m_2\end{pmatrix}.
  \]

  If \(N=2q\), applying this to all \(q\) two-row pairs of \(OT\), and placing
  the obtained \(m_1,m_2\) in the same positions, gives
  \[
    OT=HM,\qquad H^{\mathsf T}H=pI_N,
  \]
  with \(M\in\Mat_{N\times n}(\Poly)\), where \(H\in\Mat_N(\Poly)\) is
  block diagonal with \(q\) copies of \(B\).

  If \(N=2q+1\), the last row of \(OT(z)\) is zero, so every entry in that row
  is divisible by \(p\).  Together with the same \(q\) two-row blocks, this gives
  \(OT=HM\), where
  \[
    H=\operatorname{diag}(\underbrace{B,\ldots,B}_{q\text{ copies}},p)
  \]
  and \(M\in\Mat_{N\times n}(\Poly)\).  This \(H\) satisfies
  \[
    H^{\mathsf T}H=\operatorname{diag}(pI_{2q},p^2).
  \]
\end{proof}

\begin{lemma}[Even-row one-step descent for a no-real-root quadratic]
\label{lem:no-real-root-even-row-one-step-en}
  Let \(q,n\in\mathbb{N}\).
  Let \(a,b\in\mathbb{R}\) with \(b\ne0\), and put
  \(p=(X-a)^2+b^2\in\Poly\).
  If \(A\in\Mat_n(\Poly)\) and \(T\in\Mat_{2q\times n}(\Poly)\) satisfy
  \[
    pA=T^{\mathsf T}T,
  \]
  then there is \(U\in\Mat_{2q\times n}(\Poly)\) such that
  \[
    A=U^{\mathsf T}U.
  \]
\end{lemma}

\begin{proof}
  Apply the even-row case of Lemma~\ref{lem:two-row-block-construction-en} to \(T\).
  We get a real orthogonal matrix \(O\in\Mat_{2q}(\mathbb{R})\), a polynomial matrix
  \(U\in\Mat_{2q\times n}(\Poly)\), and \(H\in\Mat_{2q}(\Poly)\) such that
  \[
    OT=HU,\qquad H^{\mathsf T}H=pI_{2q}.
  \]
  Since \(O\) is orthogonal,
  \[
    pA=T^{\mathsf T}T
      =(OT)^{\mathsf T}(OT)
      =U^{\mathsf T}H^{\mathsf T}HU
      =pU^{\mathsf T}U.
  \]
  Cancelling \(p\) gives \(A=U^{\mathsf T}U\).
\end{proof}

\begin{lemma}[No-real-root quadratic factor descent]
\label{lem:no-real-root-quadratic-factor-descent-en}
  Let \(n\in\mathbb{N}\).
  Let \(p\in\Poly\) be an irreducible quadratic polynomial with no real root.
  If \(A\in\Mat_n(\Poly)\) and \(T\in\Mat_n(\Poly)\) satisfy
  \[
    p^2A=T^{\mathsf T}T,
  \]
  then there is \(U\in\Mat_n(\Poly)\) such that
  \[
    A=U^{\mathsf T}U.
  \]
\end{lemma}

\begin{proof}
  Completing the square gives \(p=c p_0\), where \(c\in\mathbb R^\times\)
  and \(p_0=(X-a)^2+b^2\), \(b\ne0\).  Replacing \(T\) by \(T/c\) turns
  \(p^2A=T^{\mathsf T}T\) into \(p_0^2A=(T/c)^{\mathsf T}(T/c)\).
  Thus it suffices to treat \(p=(X-a)^2+b^2\).

  If \(n\) is even, write \(n=2q\).  Applying
  Lemma~\ref{lem:no-real-root-even-row-one-step-en} to \(pA\) and \(T\) gives
  \(pA=M^{\mathsf T}M\) for some \(M\in\Mat_n(\Poly)\).  Applying the same
  lemma once more to \(A\) and \(M\) gives \(A=U^{\mathsf T}U\).

  Now suppose that \(n\) is odd, say \(n=2q+1\).
  Apply the odd-row case of Lemma~\ref{lem:two-row-block-construction-en} to \(T\).
  We get a real orthogonal matrix \(O\in\Mat_n(\mathbb{R})\),
  \(M\in\Mat_n(\Poly)\), and \(H\in\Mat_n(\Poly)\) such that
  \[
    OT=HM,\qquad H^{\mathsf T}H=\operatorname{diag}(pI_{2q},p^2).
  \]
  Write \(M_{\mathrm{pair}}\)
  for the first \(2q\) rows of \(M\), and \(M_{\mathrm{rem}}\) for the last row.
  Using the orthogonality of \(O\), the hypothesis becomes
  \[
    p^2A
      =(OT)^{\mathsf T}(OT)
      =
    p\,M_{\mathrm{pair}}^{\mathsf T}M_{\mathrm{pair}}
    +p^2\,M_{\mathrm{rem}}^{\mathsf T}M_{\mathrm{rem}}.
  \]
  Hence
  \[
    p\bigl(A-M_{\mathrm{rem}}^{\mathsf T}M_{\mathrm{rem}}\bigr)
      =
    M_{\mathrm{pair}}^{\mathsf T}M_{\mathrm{pair}}.
  \]
  Apply Lemma~\ref{lem:no-real-root-even-row-one-step-en} to
  \(A-M_{\mathrm{rem}}^{\mathsf T}M_{\mathrm{rem}}\) and \(M_{\mathrm{pair}}\).
  This gives
  \[
    A-M_{\mathrm{rem}}^{\mathsf T}M_{\mathrm{rem}}
      =
    U_{\mathrm{pair}}^{\mathsf T}U_{\mathrm{pair}}
  \]
  for some \(U_{\mathrm{pair}}\in\Mat_{2q\times n}(\Poly)\).  Stacking
  \(M_{\mathrm{rem}}\) below \(U_{\mathrm{pair}}\) gives
  \(U\in\Mat_n(\Poly)\) with \(A=U^{\mathsf T}U\).
\end{proof}

\begin{lemma}[Denominator descent]\label{lem:denominator-descent-en}
  Let \(n\in\mathbb{N}\).
  Suppose \(A\in\Mat_n(\Poly)\) is positive semidefinite at every real point
  and
  \[
    q^2A=T^{\mathsf T}T
  \]
  for some nonzero polynomial \(q\) and
  \(T\in\Mat_n(\Poly)\). Then there exists
  \(U\in\Mat_n(\Poly)\) such that
  \[
    A=U^{\mathsf T}U
  \]
\end{lemma}

\begin{proof}
  Choose such a representation with \(\deg q\) minimal. We show that this
  minimal \(q\) must be constant. If \(q\) is not constant,
  choose an irreducible factor \(p\) and write
  \(q=pq_0\).  Rewrite the equality as
  \[
    p^2(q_0^2A)=T^{\mathsf T}T.
  \]
  If \(p\) has a real root, apply
  Lemma~\ref{lem:real-root-factor-descent-en} to \(q_0^2A\).  If
  \(p\) has no real root, then by the classification of irreducible real
  polynomials \(p\) is quadratic, and we apply
  Lemma~\ref{lem:no-real-root-quadratic-factor-descent-en} to
  \(q_0^2A\).  In either case we get a polynomial matrix \(U\) such
  that
  \[
    q_0^2A=U^{\mathsf T}U.
  \]
  Since \(\deg q_0<\deg(pq_0)=\deg q\), this contradicts the
  minimality of \(q\). Hence \(q\) was constant. Since \(q\ne0\),
  \(U=q^{-1}T\) is polynomial and \(A=U^{\mathsf T}U\).
\end{proof}

{\renewcommand{\restateprefix}{Restated; }\fullLineMatrixSOS*}

\begin{proof}
  First prove existence of a polynomial factor, without a degree bound,
  by induction on the matrix size.  The case \(m=0\) is given by the empty
  \(1\times0\) matrix.  Assume \(m>0\).  If \(\det P=0\), regard \(P\) as a linear map
  \(\Poly^m\to\Poly^m\). Since \(\Poly\) is a PID, Smith normal form applies to
  this map: there are unimodular matrices \(C,D\in\Mat_m(\Poly)\) such that
  \[
    DPC=\operatorname{diag}(s_1,\ldots,s_{m-1},0).
  \]
  Here unimodular means invertible over \(\Poly\), so the inverses also have
  polynomial entries.  The matrix \(C\) is the change of basis in the domain,
  and \(D\) is the change of basis in the codomain.  Since the last column of
  the diagonal matrix is zero, if \(e_m\) is the last standard basis vector then
  \[
    DPCe_m=0.
  \]
  Because \(D\) is invertible over \(\Poly\), this gives \(PCe_m=0\).  Thus the
  last column \(v=Ce_m\) of \(C\) lies in the kernel of \(P\).

  From this point on we use only the domain change of basis \(C\) and transform
  the quadratic form by congruence:
  \[
    B=C^{\mathsf T}PC.
  \]
  Since \(PC\) has zero last column and \(P\) is symmetric, the last row and
  last column of \(B\) are zero. Thus
  \[
    B=
    \begin{pmatrix}
      B'&0\\
      0&0
    \end{pmatrix},
  \]
  where \(B'\in\Mat_{m-1}(\Poly)\) is again positive semidefinite at every real
  point. By induction,
  \[
    B'=R'^{\mathsf T}R',
    \qquad R'\in\Mat_{m\times(m-1)}(\Poly).
  \]
  Insert a zero last column to form
  \(\bar R\in\Mat_{m\times m}(\Poly)\). Then
  \[
    B=\bar R^{\mathsf T}\bar R.
  \]
  Therefore
  \[
    P=(C^{-1})^{\mathsf T}BC^{-1}
      =(\bar R C^{-1})^{\mathsf T}(\bar R C^{-1}).
  \]
  Finally add one zero row to \(\bar R C^{-1}\) to obtain the desired
  \((m+1)\times m\) factor for \(P\).

  If \(\det P\ne0\), consider
  \[
    \widetilde P=
    \begin{pmatrix}
      P&0\\
      0&\det P
    \end{pmatrix}.
  \]
  It is positive semidefinite on the real line and has square determinant:
  \[
    \det\widetilde P=(\det P)^2.
  \]
  The square-determinant result first gives a factorization
  \(\widetilde P=S^{\mathsf T}S\) over \(\Frac\).  Clearing common
  denominators gives \(q^2\widetilde P=T^{\mathsf T}T\) with \(T\) polynomial.
  Applying denominator descent to \(\widetilde P\), we obtain a
  polynomial matrix \(U\in\Mat_{(m+1)\times(m+1)}(\Poly)\) such that
  \[
    \widetilde P=U^{\mathsf T}U.
  \]
  Split \(U\) into its first \(m\) columns and its last column:
  \[
    U=(R\mid c),
    \qquad
    R\in\Mat_{(m+1)\times m}(\Poly),\quad
    c\in\Mat_{(m+1)\times1}(\Poly).
  \]
  Then
  \[
    U^{\mathsf T}U=
    \begin{pmatrix}
      R^{\mathsf T}R & R^{\mathsf T}c\\
      c^{\mathsf T}R & c^{\mathsf T}c
    \end{pmatrix}.
  \]
    The upper-left block of \(\widetilde P\) is \(P\), hence \(P=R^{\mathsf T}R\).

  Finally compare degrees.  For each column \(j\),
  \[
    P_{jj}=\sum_\alpha R_{\alpha j}^2.
  \]
  Over the reals, the leading coefficients of nonzero squares are positive, so
  the top-degree terms in the right hand side cannot cancel.  If the right hand
  side is zero, then every \(R_{\alpha j}\) is zero and the bound is trivial.
  Otherwise
  \[
    \deg P_{jj}=\max_\alpha 2\deg R_{\alpha j}.
  \]
  Since \(\deg P_{jj}\le2d\), every entry in the \(j\)-th column of \(R\) has
  degree at most \(d\).

  This proves the bounded SOS factorization from positive semidefiniteness.
  A bounded SOS factorization gives a Gram representation by expanding the
  factor \(R\) in the monomial basis of Definition~\ref{def:sos-gram-en} and
  taking the coefficient Gram matrix \(C^{\mathsf T}C\). Conversely, if
  \(P=Q_{d,m}^{\mathsf T}YQ_{d,m}\) with \(Y\) positive semidefinite, then for
  every real \(x\) and vector \(v\),
  \[
    v^{\mathsf T}Q_{d,m}(x)^{\mathsf T}YQ_{d,m}(x)v
    =(Q_{d,m}(x)v)^{\mathsf T}Y(Q_{d,m}(x)v)\ge0.
  \]
  Thus the Gram representation implies positive semidefiniteness on the real
  line.
\end{proof}

\section{Half-Line Corollary}

{\renewcommand{\restateprefix}{Restated; }\halfLineMarkovLukacs*}

\begin{proof}
After translating \(X\mapsto X+a\), it suffices first to treat \(a=0\).  If
\(M(x)\psd\) for all \(x\ge0\), then \(P(t)=M(t^2)\) is positive semidefinite
for every real \(t\). The full-line theorem gives
\[
  M(t^2)=R(t)^{\mathsf T}R(t).
\]
Here \(R\) has \(m+1\) rows.
Write each entry of \(R\) as an even part plus \(t\) times an even part:
\[
  R(t)=A(t^2)+tB(t^2).
\]
Since \(M(t^2)\) is even in \(t\), the odd cross term vanishes, and
\[
  M(X)=A(X)^{\mathsf T}A(X)+X B(X)^{\mathsf T}B(X).
\]
The matrices \(A\) and \(B\) are obtained from the even and odd parts of the
same rows of \(R\), so both still have \(m+1\) rows.  This proves the existence
of the required half-line SOS certificate.  Since the entries of \(P(t)=M(t^2)\) have degree at
most \(2d\), the full-line theorem gives \(R(t)\) with entry degree at most
\(d\).  The even-odd decomposition then gives entry-degree bounds \(r(d)\) for
\(A\) and \(s(d)\) for \(B\).

The Gram form is obtained from this bounded SOS representation.  Expanding the
entries of \(A\) and \(B\) in the monomial bases gives
\[
  A=CQ_{r(d),m},\qquad B=DQ_{s(d),m}.
\]
Therefore \(A^{\mathsf T}A\) and \(B^{\mathsf T}B\) are represented by the
Gram matrices \(C^{\mathsf T}C\) and \(D^{\mathsf T}D\), which are positive
semidefinite.

Returning from the translation gives the general \(a\) statement, with weight
\(X-a\).  Conversely, a Gram certificate is positive semidefinite on \(x\ge a\)
because \(Y_1,Y_2\psd\), sums of positive semidefinite matrices remain positive
semidefinite, and the scalar weight \(x-a\) is nonnegative on the half-line.
Hence positive semidefiniteness, the SOS certificate, and the Gram certificate
are equivalent.
\end{proof}

\section{Finite-Interval Corollary}

{\renewcommand{\restateprefix}{Restated; }\finiteIntervalMarkovLukacs*}

\begin{proof}
For a finite interval \([a,b]\), first pull back to the unit interval:
\[
  N(t)=M(a+(b-a)t).
\]
The equivalence between positivity of \(M\) on \([a,b]\) and positivity of
\(N\) on \([0,1]\) is the affine reindexing theorem.

For \(N\) on \([0,1]\), use the half-line chart
\[
  t=\frac{u}{1+u},\qquad u\ge0.
\]
This chart identifies the half-line with \(0\le t<1\).  The endpoint \(t=1\)
is not hit directly, but the identities obtained below are polynomial
identities, so equality on \(0\le t<1\) extends to the endpoint as well.
If the entries of \(N\) have degree at most \(2d\), then
\[
  H(u)=(1+u)^{2d}N\!\left(\frac{u}{1+u}\right)
\]
is a polynomial matrix of entry degree at most \(2d\), and it is positive
semidefinite on the half-line.  The half-line Markov--Lukacs theorem gives
\[
  H(u)=A(u)^T A(u)+u\,B(u)^T B(u).
\]
Here the entries of \(A\) have degree at most \(d\), and the entries of \(B\)
have degree at most \(\max(d-1,0)\).

  If \(d>0\), dehomogenizing by \(u=t/(1-t)\) gives the even-bound certificate
  \[
    N(t)=\widetilde A(t)^T\widetilde A(t)
      +t(1-t)\widetilde B(t)^T\widetilde B(t),
  \]
  where
  \[
    \widetilde A(t)=(1-t)^dA\!\left(\frac{t}{1-t}\right),\qquad
    \widetilde B(t)=(1-t)^{d-1}B\!\left(\frac{t}{1-t}\right).
  \]
  The degree bounds on \(A\) and \(B\) make \(\widetilde A\) and
  \(\widetilde B\) polynomial matrices of degrees at most \(d\) and \(d-1\),
  respectively.  If \(d=0\), then \(B=0\) for degree reasons and the same form
  holds without the second term.
  The weight \(t(1-t)\) is what becomes \((X-a)(b-X)\) after returning to the
  original interval.
  For an odd bound \(2d+1\), the same chart gives
\[
  H(u)=(1+u)^{2d+1}N\!\left(\frac{u}{1+u}\right),
\]
which has entry degree at most \(2d+1\) and is positive semidefinite on the
half-line.  The half-line theorem gives
\[
  H(u)=A(u)^TA(u)+uB(u)^TB(u).
\]
  The entries of \(B\) have degree at most \(d\).  The apparent degree \(d+1\)
  term in \(A\) vanishes as follows.  For each column \(j\), the diagonal entry
  satisfies
  \[
    H_{jj}=\sum_\alpha A_{\alpha j}^2
      +u\sum_\alpha B_{\alpha j}^2.
  \]
  The second term has degree at most \(2d+1\), and \(H_{jj}\) also has degree at
  most \(2d+1\).  Thus the coefficient of \(u^{2d+2}\) is the sum of the
  squares of the \(u^{d+1}\)-coefficients of the \(A_{\alpha j}\), and this sum
  is zero.  Over \(\mathbb{R}\), a sum of squares is zero only when every term
  is zero, so every \(A_{\alpha j}\) has degree at most \(d\).
  Dehomogenization then gives
  \[
    N(t)=t\,\widetilde B(t)^T\widetilde B(t)
      +(1-t)\,\widetilde A(t)^T\widetilde A(t),
  \]
  with both dehomogenized factors of degree at most \(d\).
Finally, the affine inverse \(t=(X-a)/(b-a)\) turns this into
\[
  M(X)=T_0(X)+(X-a)(b-X)T_1(X),
\]
in the even case, and
\[
  M(X)=(X-a)T_0(X)+(b-X)T_1(X)
\]
in the odd case, after absorbing the positive scalar factors into the SOS
matrices.
The Gram versions follow by expanding these bounded SOS factors in the
monomial basis and collecting coefficients into positive semidefinite Gram
matrices.
Conversely, a Gram certificate is positive semidefinite on the
interval because the Gram matrices are positive semidefinite and the weights
\(X-a\), \(b-X\), and \((X-a)(b-X)\) are nonnegative on \([a,b]\).
Thus, in each parity case, positive semidefiniteness, the SOS certificate, and
the Gram certificate are equivalent.
\end{proof}

\appendix

\section{Lean Correspondence}

The proof above corresponds to the following Lean theorems.  Names below
are relative to the namespace \lean{MatrixSOS}.  The standalone Tsen theorem
requires \lean{import MatrixSOS.Proof.Tsen}; the main certificate API is
available through \lean{import MatrixSOS}.
\begin{itemize}
  \item The auxiliary ordered Hahn field is \lean{RationalHahnField}, and its
    square-root source is \lean{rationalHahnNonnegSquareStatement}.
  \item The binary sum-of-squares representation is proved as
    \lean{fracPoly_binaryForm_represents_one_of_nonnegWhereDefined}.
  \item The common-zero lemma over an algebraically closed field is public as
    \lean{algebraicallyClosed_commonZero_of_homogeneous}.
  \item The \(K(X)\) case of Tsen's theorem, namely that \(K(X)\) is \(C_1\)
    for algebraically closed \(K\), is public as
    \lean{Tsen.algebraicallyClosed_ratFunc_c1FieldStatement}.
  \item The full-line theorem is
    \lean{fullLine_posSemidef_iff_sos}.
  \item The bounded right half-line SOS theorem is
    \lean{Certificates.HalfLine.posSemidefOn_Ici_iff_sos}, and its
    Gram form is \lean{Certificates.HalfLine.posSemidefOn_Ici_iff_gram}.
  \item The finite-interval Markov--Lukacs SOS certificates are
    \lean{Certificates.Interval.posSemidefOn_Icc_iff_lukacsSOS_even}
    and \lean{Certificates.Interval.posSemidefOn_Icc_iff_lukacsSOS_odd}.
  \item Their bounded Gram forms are
    \lean{Certificates.Interval.posSemidefOn_Icc_iff_lukacsGram_even}
    and \lean{Certificates.Interval.posSemidefOn_Icc_iff_lukacsGram_odd}.
  \item The scalar and constant special cases are
    \lean{Certificates.Scalar.polynomial_nonnegOn_univ_iff_sos} and
    \lean{Certificates.Constant.posSemidef_iff_exists_gram}.
\end{itemize}

\section{Preparation for the proof of Lemma~\ref{lem:binary-sum-two-squares-en}}

The rest of the appendix proves Lemma~\ref{lem:binary-sum-two-squares-en}.
{\renewcommand{\restateprefix}{Restated; }\binarySumTwoSquares*}
Here \(\NNwd(r)\) means that, for a representation \(r=f/g\) with \(g\ne0\),
\[
  g(x)\ne0 \quad\Longrightarrow\quad \frac{f(x)}{g(x)}\ge0
  \qquad (x\in\mathbb{R}).
\]

The proof of Lemma~\ref{lem:binary-sum-two-squares-en} proceeds as follows.
First, the rational-exponent Hahn field is used to detect a real negative value
inside a nonnegative-square-closed ordered field.  An ordered field \(K\) is
nonnegative-square-closed if every nonnegative element is a square:
\[
  0\le u \quad\Longrightarrow\quad \exists v\in K,\ u=v^2.
\]
This negative-value detection shows that a rational function which maps to a
nonnegative element in every nonnegative-square-closed ordered extension must
satisfy \(\NNwd\).  Finally, Tsen's theorem and the splitting criterion for
Pfister forms give a local-global principle for ternary quadratic forms, from
which Lemma~\ref{lem:binary-sum-two-squares-en} follows.

We use the following quadratic-form terminology in this appendix.  Over a field
\(F\), the diagonal quadratic form
\[
  \langle c_1,\ldots,c_n\rangle
\]
is the form
\[
  (x_1,\ldots,x_n)\longmapsto c_1x_1^2+\cdots+c_nx_n^2.
\]
It is isotropic if it has a nonzero zero.

\section{Rational Hahn Field}\label{sec:hahn-start-en}

We write \(\RationalHahn\) for the Hahn series field
\[
  \Hahn=\mathbb{R}((t^{\mathbb{Q}}))
\]
whose series have well-ordered support, equipped with the lexicographic
order.  For background on Hahn series, see
Hahn~\cite{Hahn1907} and Ribenboim~\cite{RibenboimFields}. If
\[
  f=\sum a_qt^q\ne0,
\]
then \(\ord(f)\) is the smallest exponent in its support and
\(\lc(f)=a_{\ord(f)}\). Positivity means \(\lc(f)>0\).

\begin{theorem}[Nonnegative elements are squares]
  If \(f\in\RationalHahn\) and \(f\ge0\), then there exists
  \(g\in\RationalHahn\) such that \(f=g^2\).
\end{theorem}

\begin{proof}
  The case \(f=0\) is immediate. Suppose \(f>0\). Let \(\alpha=\ord(f)\) and
  \(c=\lc(f)>0\). Choose \(a\in\mathbb{R}\) with \(a^2=c\). Put
  \[
    u=c^{-1}t^{-\alpha}f.
  \]
  Then \(u\) has leading term \(1\), and all terms of \(u-1\) have positive
  order. Hence the Hahn-series binomial expansion
  \[
    z=\sum_{n\ge0}\binom{1/2}{n}(u-1)^n
  \]
  is defined.  If \(u=1\), take \(z=1\).  Otherwise let
  \(\epsilon=\ord(u-1)>0\).  The \(n\)-th power has order \(n\epsilon\),
  which is unbounded in \(\mathbb Q\).  Thus below any fixed rational bound
  only finitely many powers contribute, and their supports form a finite
  union of well-ordered sets.  The total support is therefore well ordered,
  and each exponent receives contributions from only finitely many summands. The usual binomial
  coefficient calculation then gives \(z^2=u\). Therefore
  \[
    g=a\,t^{\alpha/2}z
  \]
  satisfies \(g^2=f\).
\end{proof}

\section{Detecting Negative Values}

For each real \(x\), the map \(X\mapsto x+t\) sends \(\Frac\) into the Laurent
series field around \(x\), and then into \(\RationalHahn\) by embedding integer
exponents into rational exponents.  Nonzero polynomials remain nonzero
under this substitution, so their images can be inverted; this defines a
unital field embedding and hence the required algebra structure.

\begin{lemma}\label{lem:negative-polynomial-hahn-en}
  If \(p\in\Poly\) and \(p(x)<0\) for some \(x\in\mathbb{R}\), then there is a
  \(\Frac\)-algebra structure on \(\RationalHahn\) for which the image of \(p\)
  is negative.
\end{lemma}

\begin{proof}
  Expand \(p\) around \(x\):
  \[
    p(x+T)=p(x)+a_1T+a_2T^2+\cdots.
  \]
  The constant term is negative, so this Laurent series is negative in the
  lexicographic order. The order is preserved by the embedding into rational
  Hahn series.
\end{proof}

\begin{lemma}\label{lem:nonneg-in-all-extensions-implies-nonneg-where-defined-en}
  If \(r\in\Frac\) maps to a nonnegative element in every nonnegative-square-closed ordered
  \(\Frac\)-algebra, then \(r\) is nonnegative at every real point where it is
  defined.
\end{lemma}

\begin{proof}
  If \(r=f/g\) were negative at a point \(x\) with \(g(x)\ne0\), then
  \(f(x)g(x)<0\). Lemma~\ref{lem:negative-polynomial-hahn-en} applied to the
  polynomial \(fg\) would give a nonnegative-square-closed ordered
  \(\Frac\)-algebra in which \(fg<0\), and hence
  \(r<0\), contradicting the hypothesis.
\end{proof}

\begin{lemma}[Nonnegative rational functions as sums of two squares]
\label{lem:nonneg-rational-function-sum-two-squares-en}
  If \(r\in\Frac\) is nonnegative at all real points where it is defined, then
  \[
    r=u^2+v^2
  \]
  for some \(u,v\in\Frac\).
\end{lemma}

\begin{proof}
  Write \(r=f/g\), hence \(r=fg/g^2\). The polynomial \(fg\) is nonnegative on
  the real line. A nonnegative real polynomial factors into even real roots and
  positive quadratic factors. The even-root part is a square, and each positive
  quadratic factor is a sum of two squares. The identity
  \[
    (u^2+v^2)(s^2+t^2)=(us-vt)^2+(ut+vs)^2
  \]
  shows inductively over the factors that \(fg\) is a sum of two squares.
  Dividing by \(g^2\) gives the result.
\end{proof}

\section{A Ternary Local-Global Step}\label{sec:ternary-local-global-en}

Let \(F=\Frac\). The proof uses a local-global statement for
\[
  X^2=aY^2+bZ^2.
\]
The ordered extensions used here are nonnegative-square-closed ordered fields
equipped with an \(F\)-algebra structure.

\subsection{\texorpdfstring{The \(k(X)\) Case of Tsen's Theorem}{The k(X) Case of Tsen's Theorem}}

\begin{definition}
  A field \(E\) is \(C_1\) if every homogeneous polynomial over \(E\) of
  degree \(d\ge1\) in more than \(d\) variables has a nontrivial zero over
  \(E\).
\end{definition}

The form of Tsen's theorem used here is the rational-function-field case: if \(k\) is
algebraically closed, then \(k(X)\) is \(C_1\).  A broader one-variable function
field version is also known as Tsen's theorem, but this proof only uses the
rational function field case.  The actual application is still narrower:
four-variable quadratic forms over \(\mathbb{C}(X)\).  We use the following
common-zero lemma behind the standard proof.

In this subsection, the height \(\operatorname{ht}(\mathfrak p)\) of a prime
ideal \(\mathfrak p\) is the supremum of the lengths \(n\) of chains
\[
  \mathfrak p_0\subsetneq\cdots\subsetneq\mathfrak p_n=\mathfrak p
\]
of prime ideals.  The height \(\operatorname{ht}(J)\) of an ideal \(J\) is
defined as the minimum height of a minimal prime ideal containing \(J\).
\footnote{Equivalently, it is the minimum height among all prime ideals
containing \(J\). This equivalence follows easily from lemmas available in
mathlib.}
This is the definition corresponding to Lean's \(\lean{Ideal.height}\).

\begin{theorem}[Hilbert's Nullstellensatz]\label{thm:nullstellensatz-en}
  Let \(k\) and \(K\) be fields, assume that \(K\) is algebraically closed, and
  assume that \(K\) is a \(k\)-algebra.  Put \(A=k[Z_1,\ldots,Z_M]\).  For any
  ideal \(J\subset A\), let \(V_K(J)\subset K^M\) be its \(K\)-valued common
  zero locus, and let \(I_k(V_K(J))\) be the ideal of \(k\)-coefficient
  polynomials vanishing on that set. Then
  \[
    I_k(V_K(J))=\sqrt J.
  \]
\end{theorem}

\begin{proof}
  In the formalization this is exactly the mathlib theorem
  \begin{quote}
    \lean{MvPolynomial.vanishingIdeal_zeroLocus_eq_radical}
  \end{quote}
\end{proof}

\begin{lemma}[Vanishing ideal of the origin]\label{lem:vanishing-ideal-origin-en}
  Let \(k\) be a field, and let \(A=k[Z_1,\ldots,Z_M]\). Then
  \[
    I(\{0\})=(Z_1,\ldots,Z_M).
  \]
\end{lemma}

\begin{proof}
  A polynomial \(f\in A\) vanishes at the origin if and only if its constant
  coefficient is zero.  The variable ideal \((Z_1,\ldots,Z_M)\) is exactly the
  set of polynomials with zero constant coefficient.  Hence the two ideals are
  equal.
\end{proof}

\begin{lemma}[Radicals do not change height]\label{lem:height-radical-en}
  If \(A\) is a commutative ring and \(J\subset A\) is an ideal, then
  \[
    \operatorname{ht}(J)=\operatorname{ht}(\sqrt J).
  \]
\end{lemma}

\begin{proof}
  This is formalized as \(\lean{ideal_height_radical}\).
  A prime ideal \(\mathfrak p\) contains \(J\) if and only if it contains
  \(\sqrt J\).  Therefore the minimal prime ideals containing \(J\) and those
  containing \(\sqrt J\) are the same, so the minimum of their heights is the
  same.
\end{proof}

\begin{lemma}[Height of the variable ideal]\label{lem:height-ideal-of-vars-en}
  Let \(k\) be a field, and let \(A=k[Z_1,\ldots,Z_M]\). Then
  \[
    \operatorname{ht}(Z_1,\ldots,Z_M)=M.
  \]
\end{lemma}

\begin{proof}
  This is formalized as \(\lean{idealOfVars_height_fintype}\).
  The maximality of the variable ideal is obtained directly from the fact that
  it is the kernel of the constant-coefficient map \(k[Z_1,\ldots,Z_M]\to k\).
  The empty-variable case is immediate.  For the induction step, Lean identifies
  a multivariable polynomial ring with one more variable as a polynomial ring in
  one variable over the previous multivariable polynomial ring, using
  \(\lean{MvPolynomial.optionEquivLeft}\).  Heights are transported across this
  equivalence by \(\lean{RingEquiv.height_map}\), and the height increases by
  one by \(\lean{Polynomial.height_eq_height_add_one}\).  Iterating this over a
  finite set of variables gives height \(M\).
\end{proof}

\begin{theorem}[Krull's height theorem]\label{thm:krull-height-en}
  Let \(A\) be a Noetherian commutative ring, and let \(J\subset A\) be a
  proper ideal. Then
  \[
    \operatorname{ht}(J)\le \operatorname{spanFinrank}(J).
  \]
  Here \(\operatorname{spanFinrank}(J)\) is mathlib's finite spanning rank,
  \(\lean{Submodule.spanFinrank}\).  Since ideals in a Noetherian ring are
  finitely generated, this is the minimum cardinality of a finite generating
  set of \(J\).
\end{theorem}

\begin{proof}
  Lean uses the mathlib theorem \(\lean{Ideal.height_le_spanFinrank}\).
\end{proof}

\begin{lemma}[Common-zero lemma]\label{lem:projective-common-zero-en}
  Let \(k\) be algebraically closed, and let
  \(P_1,\ldots,P_r\in k[Z_1,\ldots,Z_M]\) be homogeneous polynomials of positive
  degree.  If \(r<M\), then they have a common nontrivial zero in \(k^M\).
\end{lemma}

\begin{proof}
  Let \(I=(P_1,\ldots,P_r)\). Suppose, for contradiction, that the only common
  zero is the origin.

  In Lean, \(\lean{MvPolynomial.zeroLocus_span}\) first identifies the zero
  locus of \(I\) with the common zero locus of the generators \(P_i\).  Since
  the \(P_i\) are homogeneous of positive degree, the following mathlib lemmas
  show that the origin is always a zero:
  \begin{quote}
    \lean{MvPolynomial.IsHomogeneous.coeff_eq_zero}\\
    \lean{MvPolynomial.constantCoeff_eq}\\
    \lean{MvPolynomial.aeval_zero}
  \end{quote}

  Applying Theorem~\ref{thm:nullstellensatz-en} to \(I\), and then using
  Lemma~\ref{lem:vanishing-ideal-origin-en}, gives
  \[
    \sqrt I=(Z_1,\ldots,Z_M).
  \]
  By Lemma~\ref{lem:height-radical-en} and
  Lemma~\ref{lem:height-ideal-of-vars-en}, the height of \(I\) is \(M\).

  On the other hand, Theorem~\ref{thm:krull-height-en} applied to \(I\) gives
  \(\operatorname{ht}(I)\le\operatorname{spanFinrank}(I)\).  In the
  formalization,
  \(\lean{Submodule.spanFinrank_span_le_ncard_of_finite}\) and
  \(\lean{Fintype.card_range_le}\) bound this finite spanning rank by \(r\),
  because \(I\) is generated by the \(r\) polynomials \(P_i\). Thus
  \(M\le r\), contradicting \(r<M\).
\end{proof}

\begin{theorem}[The \(k(X)\) case of Tsen's theorem]
\label{thm:tsen-rational-function-field-en}
  If \(k\) is algebraically closed, then \(k(X)\) is \(C_1\).  In particular,
  \(\mathbb{C}(X)\) is \(C_1\).
\end{theorem}

\begin{proof}
  Let \(f(Y_1,\ldots,Y_n)\in k(X)[Y_1,\ldots,Y_n]\) be homogeneous of degree
  \(d\), with \(n>d\).  After clearing denominators, assume that the
  coefficients of \(f\) lie in \(k[X]\), and let \(D\) be the maximum \(X\)-degree
  of those coefficients.  Put \(N=D\), and substitute unknown polynomials
  \[
    Y_i(X)=c_{i0}+c_{i1}X+\cdots+c_{iN}X^N .
  \]
  After substitution, every coefficient of the resulting polynomial in \(X\)
  is a homogeneous degree-\(d\) polynomial in the \(c_{ij}\).  There are
  \(n(N+1)\) variables.  The number of coefficient equations is at most
  \(D+dN+1\), because each substituted term has \(X\)-degree at most \(D+dN\).
  Since \(n>d\), for \(N=D\) one has \(n(N+1)>D+dN+1\).  Lemma
  \ref{lem:projective-common-zero-en} therefore gives a nontrivial choice of
  the \(c_{ij}\), hence a nontrivial zero of \(f\) over \(k(X)\).
\end{proof}

\subsection{Pfister Forms and the Ternary Local-Global Step}

\begin{lemma}\label{lem:ternary-local-global-en}
  Let \(a,b\in F^\times\). If \(X^2=aY^2+bZ^2\) has a nontrivial solution over
  every nonnegative-square-closed ordered field equipped with an \(F\)-algebra
  structure, then it has a nontrivial solution over \(F\).
\end{lemma}

\begin{proof}
  Consider the associated two-fold Pfister form
  \[
    \pi=\langle 1,-a,-b,ab\rangle.
  \]
  Over \(F(i)\simeq\mathbb{C}(X)\),
  Theorem~\ref{thm:tsen-rational-function-field-en} gives isotropy of the
  four-variable quadratic form \(\pi\).  Write a nonzero isotropic vector as \(u+iv\), with \(u,v\in F^4\).
  Its real and imaginary parts give
  \(\pi(u)=\pi(v)\) and \(B_\pi(u,v)=0\), where \(B_\pi\) is the associated
  symmetric bilinear form.  If this common value is zero, one of \(u,v\)
  is already a nonzero \(F\)-rational isotropic vector.
  Otherwise the common value is \(c\ne0\), so \(u,v\) span a nonsingular
  plane with diagonal form \(\langle c,c\rangle\).  Diagonalizing its
  orthogonal complement gives \(\langle c,c,d,e\rangle\).
  Since \(\det\pi=a^2b^2\) and determinants under congruence differ by a
  nonzero square, \(de=s^2\) for some \(s\ne0\).  Multiplying the last
  basis vector by \(d/s\) changes its coefficient to \(d\).
  Thus the form is transformed into
  \[
    \langle c,c,d,d\rangle.
  \]
  The remaining diagonal entries are nonzero, since the Pfister form is
  nonsingular.
  The local isotropy hypothesis implies that this diagonal form is isotropic
  over every nonnegative-square-closed ordered field equipped with an
  \(F\)-algebra structure. For nonnegative-square-closed ordered fields this is
  equivalent to nonnegativity of \(-cd\): in a nonzero zero of
  \(c(u_1^2+u_2^2)+d(u_3^2+u_4^2)\), the first sum of squares is nonzero,
  giving \(-c/d=(u_3^2+u_4^2)/(u_1^2+u_2^2)\ge0\).
  Conversely a square root \(r^2=-cd\) gives the zero \((1,0,r/d,0)\).
  Lemma~\ref{lem:nonneg-in-all-extensions-implies-nonneg-where-defined-en}
  and Lemma~\ref{lem:nonneg-rational-function-sum-two-squares-en} show that
  \(-cd\) is a sum of two squares in \(F\), say \(-cd=r^2+s^2\). Then
  \((d,0,r,s)\) is an isotropic vector for \(\langle c,c,d,d\rangle\). It is
  nonzero because \(d\ne0\). Transporting it back gives an isotropic vector for
  \(\pi\) with coordinates in \(F\).
  The two-fold Pfister splitting criterion~\cite{LamQuadraticForms} can here
  be proved directly.  Let \((x,y,z,w)\) be such a nonzero zero and put
  \(D=z^2-aw^2\), so \(x^2-ay^2=bD\).  If \(D=0\), either \((z,w)\ne(0,0)\)
  gives \(a=(z/w)^2\), or \((x,y)\ne(0,0)\) gives \(a=(x/y)^2\).
  In either case \(a\) is a square and the ternary equation has a nonzero
  solution with \(Z=0\).  If \(D\ne0\), the identity
  \[
    (xz-a y w)^2-a(xw-yz)^2=(x^2-a y^2)(z^2-a w^2),
  \]
  gives the explicit nontrivial solution
  \[
    (X,Y,Z)=\left(\frac{xz-ayw}{D},\frac{xw-yz}{D},1\right)
  \]
  of \(X^2=aY^2+bZ^2\).  The ternary form used here is the subform
  \(\langle1,-a,-b\rangle\); the conventional pure subform of \(\pi\)
  is instead \(\langle-a,-b,ab\rangle\).
\end{proof}

\begin{lemma}[Making the first coordinate nonzero]
\label{lem:first-coordinate-nonzero-en}
  Let \(E\) be a field of characteristic \(0\), and let \(a,b\in E^\times\).
  If \(x^2=ay^2+bz^2\) has a nontrivial solution over \(E\), then it has one
  with \(x\ne0\).
\end{lemma}

\begin{proof}
  If the original solution has \(x_0\ne0\), there is nothing to prove.  If
  \(x_0=0\), then \(y_0,z_0\ne0\) and \(a y_0^2+b z_0^2=0\).  Set
  \[
    x=2ay_0,\qquad y=(a+1)y_0,\qquad z=(1-a)z_0 .
  \]
  Then \(x\ne0\), and
  \[
    a((a+1)y_0)^2+b((1-a)z_0)^2
    =(2ay_0)^2+(1-a)^2(a y_0^2+b z_0^2)
    =(2ay_0)^2.
  \]
\end{proof}

{\renewcommand{\restateprefix}{Restated; }\binarySumTwoSquares*}

\begin{proof}
  First show that
  \[
    X^2=aY^2+bZ^2
  \]
  has a nontrivial solution over every nonnegative-square-closed ordered
  \(F\)-algebra. By
  Lemma~\ref{lem:nonneg-rational-function-sum-two-squares-en}, the functions
  \(a\) and \(b\) are sums of two squares in \(F\). In any such ordered field,
  sums of squares are nonnegative and the images of \(a,b\) are nonzero, hence
  positive. Since nonnegative elements are squares, the image \(A\) of \(a\)
  has a square root \(\alpha\), and
  \[
    (X,Y,Z)=(1,\alpha^{-1},0)
  \]
  is a nontrivial solution.
  Lemma~\ref{lem:ternary-local-global-en} gives a nontrivial solution over
  \(F\), and Lemma~\ref{lem:first-coordinate-nonzero-en} lets us choose it
  with \(X\ne0\). Dividing by \(X\) gives \(u=Y/X\) and \(v=Z/X\), hence
  \[
    1=au^2+bv^2.
  \]
\end{proof}

\begin{thebibliography}{99}

\bibitem{HanselkaSinn}
Christoph Hanselka and Rainer Sinn,
\emph{Positive Semidefinite Univariate Matrix Polynomials},
Mathematische Zeitschrift 292 (2019), no. 1--2, 83--101.

\bibitem{BPSV}
Grigoriy Blekherman, Daniel Plaumann, Rainer Sinn, and Cynthia Vinzant,
\emph{Low-Rank Sum-of-Squares Representations on Varieties of Minimal Degree},
International Mathematics Research Notices 2019, no. 1, 33--54.

\bibitem{KlepSchweighofer}
Igor Klep and Markus Schweighofer,
\emph{Pure States, Positive Matrix Polynomials and Sums of Hermitian Squares},
Indiana University Mathematics Journal 59 (2010), no. 3, 857--874.

\bibitem{PowersReznick}
Victoria Powers and Bruce Reznick,
\emph{Polynomials That Are Positive on an Interval},
Transactions of the American Mathematical Society 352 (2000), no. 10,
4677--4692.

\bibitem{LamQuadraticForms}
Tsit-Yuen Lam,
\emph{Introduction to Quadratic Forms over Fields},
Graduate Studies in Mathematics, vol. 67, American Mathematical Society, 2005.

\bibitem{LangAlgebra}
Serge Lang,
\emph{Algebra},
Graduate Texts in Mathematics, vol. 211, revised third edition, Springer, 2002.

\bibitem{EisenbudCA}
David Eisenbud,
\emph{Commutative Algebra with a View Toward Algebraic Geometry},
Graduate Texts in Mathematics, vol. 150, Springer, 1995.

\bibitem{Hahn1907}
Hans Hahn,
\emph{\"Uber die nichtarchimedischen Gr\"ossensysteme},
Sitzungsberichte der Kaiserlichen Akademie der Wissenschaften, Wien,
Mathematisch-Naturwissenschaftliche Klasse 116 (1907), 601--655.

\bibitem{RibenboimFields}
Paulo Ribenboim,
\emph{The Theory of Classical Valuations},
Springer Monographs in Mathematics, Springer, 1999.

\end{thebibliography}

\end{document}
